\chapter{Interpretation and Its Limits}
\label{ch:philosophy}

What do the preceding chapters commit to, and which readings of them could an
observation or a proof count against? The chapter's readings are definitions,
hypotheses, and open problems, and it proves a proposition and a theorem.
\Cref{prop:phil-seeded-contrast-formation} shows that on the four-agent cycle
of \cref{sec:rg-formation-endogenous-cycle} a collective forms from every
Gaussian initial law with probability bounded below uniformly in the regime
parameter, whether or not the source prior learns, because the random initial
assignment seeds the pair; its proof is in \cref{sec:rg-formation-cycle-proof}.
The main statement, \cref{conj:phil-contrast-formation}, sharpens the
participatory reading: at a sufficiently small fixed temperature, slow learning of the source
prior selects one pairing at the level of the law from almost every Gaussian
initial law, while a frozen prior leaves the long-time pairing law symmetric.
It feeds \cref{open:phil-testability}.

Beyond \cref{sec:intro-conventions}, the chapter declines the claims that a
structural-realist, Neo-Kantian, or noumenal reading follows from the invariance
results or is selected by the mathematics; that the lineage citations below
prove the probability construction or establish its interpretation; that
\cref{def:phil-idle-wheel} follows from constructive empiricism; that
participation entails consciousness or that variational free energy is a
physical action principle; and that a constructed record sensitive to link
holonomy is a canonical observable or an empirical realization. Its open items
are \cref{claim:physical-law-identification,open:phil-testability}.
\status{NOT-CLAIMED}

\section{Scaffolding, idle posits, and structural readings}
\label{sec:phil-base}
\label{sec:phil-noumenon}
\label{sec:phil-realism}

The verdict of the preceding chapters is that the architecture is local-first,
not background-free. Agent \(i\) carries the law points
\begin{equation}
q_i^b\in(\mathcal E_b)_{c_*},
\qquad
q_i^m\in(\mathcal E_m)_{c_*},
\label{eq:phil-fiber-measures}
\end{equation}
which the cross maps \(\Phi\) and \(\widetilde\Phi\) do not identify, and the
base, principal bundle, and measurable spaces are coordination scaffolds that
make agent-relative interfaces comparable and that local inference laws do not
construct. The verdict would change if edges became random structural
coordinates updated by inference (\cref{open:obs-edge-creation}) or links were
recovered from population records, with single-point part
\cref{open:phil-testability}. A Neo-Kantian reading takes local frames, model
evaluations, and access maps as constitutive but revisable conditions of a
record's inferential meaning, in the lineage of Kant, Reichenbach, Cassirer, and
Friedman \citep{Kant1781,Reichenbach1965,Cassirer1923,Friedman2001}; a noumenal
reading takes the base and bundles as structure reached only through
agent-side representations. As scaffolding is no inferential target, the
base's lack of a law decides neither reading.

\definitionheading{Idle-wheel criterion}{def:phil-idle-wheel}
A posit is idle when changing it alters the law of no record and of no
statistic of a record in the model. The idle-wheel criterion removes idle
posits by parsimony. It is a declared rule of reading.

The criterion is close in spirit to constructive empiricism, on which accepting
a theory involves belief only in its empirical adequacy
\citep[11--12]{vanFraassen1980}, but it is declared, not derived. A closed walk
\(p\) and an edge copy \(e\) determine the classes
\begin{equation}
[\mathrm{Hol}^b(p)]_{\rm conj},
\qquad
[\mathrm{Hol}^m(p)]_{\rm conj},
\qquad
[\mathring\Theta^x_e]_{\rm conj}
\label{eq:phil-holonomy}
\end{equation}
of the link holonomies \eqref{eq:geo-link-holonomy} and the per-edge twists,
invariant under passive and active changes
(\cref{prop:geo-edge-twist-invariance}). On a graph-flat family the loop classes
are trivial (\cref{prop:geo-trivializing-criterion}) and give a noumenal reading
no work; a nontrivial one is a referent only within the interaction complex,
made operational at \(c_*\) by \cref{thm:cg-loop-operational-holonomy}.

The interface theory of perception permits an arbitrary measurable relation
between representation and world \citep{Hoffman2015,Hoffman2019}. Here the
principal fiber at \(c_*\) is a \(G\)-torsor, a frame is the choice of a
reference point, and what survives a change of frame is structural data,
\begin{equation}
\mathcal I_{\mathrm{struct}}
=\left\{
\text{ELBO and relative-entropy values},\
[\mathrm{Hol}^b(p)]_{\rm conj},\
[\mathrm{Hol}^m(p)]_{\rm conj},\
\bigl[(\mathring\Theta^b_e,\mathring\Theta^m_e)_{e}\bigr]
\right\},
\label{eq:phil-invariants}
\end{equation}
the last entry being the per-edge twists of both channels modulo one
simultaneous conjugation, invariant by
\cref{prop:gen-product-evidence-invariance},
\eqref{eq:cg-kl-bimeasurable-invariance}, and
\cref{prop:geo-edge-twist-invariance}. Of the three agreements of
\cref{prop:geo-three-agreements}, frame agreement is independent of the other
two, law coherence implies mean agreement but not conversely, and none selects a
partition (\cref{open:cg-natural-selector}). These invariants are compatible
with epistemic structural realism \citep{Worrall1989}, with moderate structural
realism, on which objects share the ontological footing of relations and are
characterized only by them \citep{esfeld2008moderate}, and with ontic structural
realism, on which structure is all that exists \citep{Ladyman2007,Ladyman2014}.
No argument from theory change is given, since these are invariants of passive
frame change within one theory whereas Worrall's structure is retained across
theories, and nothing here selects among these positions or the Neo-Kantian
reading. Support would need, respectively, an invariant retained across
a change of theory, a typed interpretation of the formalism's objects and
relations in the moderate ontology, and an argument from formal
unidentifiability to ontological elimination with discriminating evidence.

\section{Participation, formation, and testability}
\label{sec:phil-participatory}
\label{sec:phil-testable}

\hypothesisheading{Operational observational closure}{hyp:phil-observational-closure}
Every datum directly available to an interior agent is a clamped state or a
message of another interior agent or of a boundary environment node, an
operational environment node in the sense of
\cref{def:obs-operational-environment-node}. The hypothesis excludes data that
reach a record kernel other than through such a node. It does not remove the
conditioning record, does not require a boundary node to be autonomous, and does
not exclude an unobserved environment that affects all agents without becoming
the argument of any agent's belief.

By \cref{thm:obs-agent-interaction-equivalence} the hypothesis is free for
kernels but removes no conditioning: a statistic of the record is
posterior-sufficient exactly when the conditional mutual information in
\eqref{eq:obs-conditioning-statistic-chain} vanishes, and a latent-independent
kernel is a zero-information boundary message
\eqref{eq:obs-conditioning-information}, not an agent.

\hypothesisheading{Every boundary node is called an agent}{hyp:phil-boundary-agents}
Every boundary environment node of \cref{hyp:phil-observational-closure} is
called an agent. The hypothesis is a naming decision. It adds to no node a
persistent state, an action, a Markov blanket, or a local free energy, and it
excludes reading the constructed joint \(\mathbb P_{\theta,V}\) or the selected
coupling \(\mathbb Q_{V,o,X}\) as an agent: these are complete joint laws at the
population level. A block conditional is a meta-agent only after its own coarse
state, observation interface, evaluated generative mechanism, recognition law,
and update rule have been constructed, as in
\cref{ch:agent-network-rg,ch:full-graph-meta-agent-vfe}. A strong agent-only
ontology would additionally need a criterion of agency for every boundary node,
and the formalism supplies no unique one.

Participation is reciprocal inferential dependence through records inside one
joint, with no live posterior in a same-step generative factor
(\cref{req:gen-typing-prohibition}). \Cref{sec:full-graph-closing-tower} locates
it in three sequences of fixed joints: the same-time inference loop, block
coordinate ascent as in any connected hierarchical model under mean-field
updates, which resembles Wheeler's participatory proposal \citep{Wheeler1990}
without being derived from it or evidence for it; parameter learning, whose
optimal prior over an unrestricted family is the aggregate recognition law
\eqref{eq:obs-aggregate-prior}; and delayed feedback through
\eqref{eq:obs-episode-update}, which for an update reading a declared
recognition kernel is the free energy of one joint
(\cref{prop:full-graph-episode-lifted-joint}) and whose dynamic closure is
part of \cref{open:rg-pointwise-parent-dynamics}. None of these separates the program
from a hierarchical model with learned priors; formation could.
\Cref{thm:finite-depth-adaptive-formation} forms successive collectives under a
supplied energy hierarchy and source support, without showing that source
learning creates the contrasts, and \cref{prop:rg-formation-endogenous-cycle}
forms one collective on a homogeneous four-cycle with probability bounded below
uniformly in its regime parameter, from a symmetric preparation whose random
initial source assignment seeds the pair, and without parameter learning. The
same seed forms a collective from every Gaussian initial law, whether or not the
source prior learns.

\propositionheading{Seeded contrast formation with idle source learning}{prop:phil-seeded-contrast-formation}
Take the homogeneous four-cycle experiment of
\cref{def:rg-formation-cycle-experiment}, with its regime parameter
\(0<\epsilon\le1/100\), the constants of \cref{eq:rg-formation-cycle-regime},
and its records, evaluator, reference law, and kinetics, and modify it in two
ways. First, the uniform prior on assignments becomes a product
\(\pi_S=\prod_i\pi_{i,S_i}\) of receiver rows, each a generative parameter that
starts uniform and learns from the recognition law \(q_t\) of the kinetics by
\begin{equation}
\dot\pi_{ij}=\epsilon\bigl(q_t(S_i=j)-\pi_{ij}\bigr),
\label{eq:phil-source-learning}
\end{equation}
the natural-gradient flow of the free energy in the prior row under the Fisher
metric of the simplex at mobility \(\epsilon\); its rest point is the recognition
label marginal, as for the aggregate prior \eqref{eq:obs-aggregate-prior}, and the
assignment is redrawn from posterior rows computed with the current prior.
Second, the initial law is \(\mathcal N(\mathsf m_0,I_{12})\) on configurations,
with an arbitrary mean \(\mathsf m_0\in\R^{12}\), times the uniform law on
assignments. Call contrast formation the event that some connected block \(I\)
with \(2\le|I|\le|V|-1\) has, at time zero, an internal edge whose belief
transported relative entropy exceeds \(1/4\), and, on some time interval of
length one inside \([0,2/\epsilon]\), all its internal belief and model
transported relative entropies at most \(\epsilon^4\) while every belief
transported relative entropy on an edge joining it to its complement exceeds
\(1/2\), the bounds of part (a) of \cref{prop:rg-formation-endogenous-cycle};
the first clause makes the event an entrance into the contrast and excludes a
contrast present at time zero that persists. Then:

(a) \emph{Rows.} The kinetics and \cref{eq:phil-source-learning} have exactly
one joint solution on \([0,\infty)\), with continuous deterministic rows, and
\(|\pi_{ij}(t)-1/2|\le\epsilon t/2\) for all \(t\ge0\).

(b) \emph{Formation.} For every \(\mathsf m_0\in\R^{12}\) and every
\(\epsilon\in(0,1/100]\),
\[
\begin{aligned}
\Pr(\text{contrast formation})
&\ge\frac{e^{-1/4}}{16}\sum_{\pm}\mathcal N(\mathsf m_0,I_{12})
\bigl\{|\mathbf x\mp\mathbf x^\circ|<\varrho_\circ\bigr\}\\
&\ge\frac{e^{-1/4}\varrho_\circ^{12}}{737280}
\exp\Bigl[-\tfrac12\bigl(|\mathbf x^\circ-\mathsf m_0|+\varrho_\circ\bigr)^2\Bigr]>0,
\end{aligned}
\]
and the same bound holds when the rows are frozen at \(1/2\). The bound is
carried by the sub-event in which \(I\) is a matched pair of the initial
assignment, held throughout \([0,t^\circ_*]\), and the contrast is that of
part (a) of \cref{prop:rg-formation-endogenous-cycle}. This sub-event is a
coherence episode with the restoring interaction and rates of that part, and
for every fixed row configuration \(\pi\) the kinetics at rows \(\pi\) leave
\(q^\star_\pi(\mathbf x,S)\propto\pi_Se^{-U_S(\mathbf x)/\mathrm T}\) invariant
and, started from it, realize the sub-event with probability at most
\(4e^{-1/(4\mathrm T)}\); so for every \(\mathsf m_0\), in either arm, the
sub-event exceeds every such stationary value by a margin bounded below as
\(\epsilon\to0\), the first display not depending on \(\epsilon\).

(c) \emph{Invariant initial laws.} The initial law is invariant under the cycle
automorphisms, with link signs transformed under reflections, exactly when
\(\mathsf m_0=(\mathrm c_1\mathbf 1_4,\mathrm c_2\mathbf 1_4,0)\) for some
\(\mathrm c_1,\mathrm c_2\in\R\), and then \(\pi_{ij}(t)=1/2\) for every
receiver, source, and \(t\ge0\).
\status{ESTABLISHED}

The proof is in \cref{sec:rg-formation-cycle-proof}. The probability in part (b)
is supplied by the preparation's matching seed, an initial assignment drawn
uniformly and independently of a configuration law of full support, and source
learning is idle in it. A statement in which learning creates a contrast must
therefore compare learned and frozen priors and concern what the law selects,
not the probability of a sample-path episode.

\theoremheading{Pairing selection by slow source learning}{conj:phil-contrast-formation}
Take the homogeneous four-cycle experiment of
\cref{def:rg-formation-cycle-experiment}, with its records, evaluator, reference
law, and kinetics, at a fixed temperature \(\mathrm T>0\) in place of
\(\epsilon^8\); let the product receiver rows of
\cref{prop:phil-seeded-contrast-formation} start uniform and learn by
\cref{eq:phil-source-learning} with a mobility \(\epsilon_{\mathrm L}>0\) in
place of \(\epsilon\); and let the initial law be
\(\mathcal N(\mathsf m_0,I_{12})\) on configurations times the uniform law on
assignments. Write \(q^{\epsilon_{\mathrm L}}_t\) for the recognition law of the
kinetics and \(q^0_t\) for that of the frozen-prior arm, whose rows stay at
\(1/2\). First, there is \(\mathrm T_0>0\) such
that for every \(\mathrm T\in(0,\mathrm T_0]\) there is
\(\mathrm C_{\mathrm T}<\infty\) for which, for Lebesgue-almost every
\(\mathsf m_0\in\R^{12}\), one of the two perfect matchings
\(M\in\{S^{(0)},S^{(1)}\}\) satisfies
\[
\lim_{\epsilon_{\mathrm L}\to0}\ \sup_{t_{\mathrm L}\le t\le2t_{\mathrm L}}
q^{\epsilon_{\mathrm L}}_t(S\ne M)=0,
\qquad
t_{\mathrm L}=\mathrm C_{\mathrm T}\epsilon_{\mathrm L}^{-1}\log(1/\epsilon_{\mathrm L}).
\]
Second, for every \(\mathrm T>0\) and every \(\mathsf m_0\),
\(q^0_t(S=S^{(0)})-q^0_t(S=S^{(1)})\to0\) as \(t\to\infty\).
\status{ESTABLISHED}

The proof is in \cref{sec:phil-pairing-selection-proof}.
\corollaryheading{Joint selection along a temperature schedule}{cor:phil-joint-selection}
In the model of \cref{conj:phil-contrast-formation}, fix any sequence
\(\mathrm T_n\downarrow0\) within its temperature interval and any
\(\eta_n\in(0,1)\) tending to zero. There are thresholds \(d_n\downarrow0\), independent
of the initial mean, such that for Lebesgue-almost every fixed
\(\mathsf m_0\), all sufficiently large \(n\), and every
\(0<\epsilon_{\mathrm L}\le d_n\), a matching \(M_n(\mathsf m_0)\) satisfies
\[
 \sup_{t_{\mathrm L}\le t\le2t_{\mathrm L}}
 q_t^{\epsilon_{\mathrm L},\mathrm T_n}(S\ne M_n)\le\eta_n,
 \qquad t_{\mathrm L}=2\epsilon_{\mathrm L}^{-1}\log(1/\epsilon_{\mathrm L}).
\]
The same thresholds can make the learned law's relative entropy to its
instantaneous invariant law, and the frozen law's relative entropy to the
uniform-row invariant law, at most \(\eta_n\) throughout this window.
Consequently the frozen matching probabilities tend to \(1/2\) each.
For any prescribed \(a>7/2\), the thresholds can also satisfy
\(d_n\le\mathrm T_n^a\). They are existence thresholds, rather than a
proved power-law selection condition; the selected matching may depend on
\(n\).
\status{ESTABLISHED}

The proof and an explicit temperature-dependent tracking bound are in
\cref{sec:phil-joint-selection-proof}. Setting
\(\mathrm T(\epsilon_{\mathrm L})=\mathrm T_n\) for
\(d_{n+1}<\epsilon_{\mathrm L}\le d_n\) gives a joint limit with a schedule
common to almost every initial mean.

The limit \(\epsilon_{\mathrm L}\to0\) at fixed \(\mathrm T\) places the
mobility of the rows below any positive mixing rate of the kinetics at fixed
rows, and \(t_{\mathrm L}\) is, up to its constant, the time in which an imprint
of order \(\epsilon_{\mathrm L}\) on the rows that grows at a rate of order
\(\epsilon_{\mathrm L}\) reaches order one. The asymmetry that learning would
amplify is supplied by the initial law, not by the generative model or the
kinetics: by part (c) of \cref{prop:phil-seeded-contrast-formation}, learning
from a Gaussian initial law that is invariant under the cycle automorphisms
leaves the rows uniform for all time. The theorem amplifies the asymmetry of
\(\mathcal N(\mathsf m_0,I_{12})\) and excludes a null set of means. The frozen
arm is the null against which selection is measured. This is a result about
the one-time assignment law at fixed temperature. The conditional block
events, a power-law vanishing-temperature regime, and conditional and dynamical parent
errors of \cref{open:rg-formation-adaptive-family} remain separate obligations.
The participatory reading requires collectives constituted by the coupled
inference, with a restoring interaction and a dependence on the preceding
state. Assignment selection alone does not establish that recursive claim. The noumenal reading buys structure behind the
agents' representations with posits that \cref{def:phil-idle-wheel} removes
unless an observable such as a loop holonomy carries them.

\openproblemheading{Physical-law identification}{claim:physical-law-identification}
Identifying an object produced by the coarse-graining of
\cref{ch:agent-network-rg}, or by its limits, with a physical law requires a
named target system, an observable and an estimator, a validation and
uncertainty protocol set before the data, a baseline margin, probability-level
closure, and held-out evidence that is robust across admissible blocking
schemes. None is supplied here, and the mathematics neither proves a physical
law nor supports ontological closure. For a named target system whose data the
model does not generate, observed through a declared record kernel that, under
the model's own law, separates two objects of the coarse-graining that differ
in a gauge-invariant datum, with an observable, an estimator, and a validation
and uncertainty protocol set before the data, the problem is settled by
held-out evidence, robust across admissible blocking schemes, that an object of
the coarse-graining meets a baseline margin with probability-level closure, or
by held-out evidence under the same protocol that the baseline margin is not
met.

Graph flatness (\cref{prop:geo-trivializing-criterion}) and the fixed-sector
condition of \cref{prop:cg-kernel-holonomy} check declared data but detect no
bundle topology or base curvature, so an interpretation must keep graph-link
holonomy, base-connection holonomy, bundle topology, and gauge redundancy apart.

\openproblemheading{Single-point testability}{open:phil-testability}
State an observation at \(c_*\) that could count against a claim of the program.
What holds is the following: \cref{thm:cg-loop-operational-holonomy} constructs a
population record and a frame-invariant statistic whose law separates two
independent-link configurations that differ only in the conjugacy class of a
loop holonomy, and \cref{conj:phil-contrast-formation} establishes a learned
pairing selection against a symmetric frozen-prior arm in one declared
family. What is missing is a
target system: sampling units, a noise model, an estimator of gauge-typed
parameters with its uncertainty, a held-out statistic, an acceptance tolerance
set in advance, and baselines without link holonomy or without formation. A
cross-scale test additionally needs two or more admissible partitions,
coarse channels, gauge alignment, and held-out coarse statistics predicted from
the fine fit without coarse refitting, and its failure would falsify the added
cross-scale hypothesis, not the finite-dimensional theorems. For a named
target system at \(c_*\) whose data the model does not generate, observed
through a declared record kernel, with sampling units and a noise model, that
reads the links of a closed walk or the coordinates of a candidate block, the
problem is settled by a test fixed before the data, with an estimator of
gauge-typed parameters, a held-out statistic, an acceptance tolerance, and
baselines without link holonomy or without formation, whose failure would
falsify the presence of nontrivial holonomy on that walk or of formation of
that block by the inference in the target, together with a proof that under
the declared kernel the law of the tested statistic depends on that
hypothesis; or, negatively, by a proof that under the declared kernel every
statistic of the records has the same law for two link configurations that
differ only in the conjugacy class of the walk's holonomy, or the same law
whether or not the block is coherent at the declared resolution.

\paragraph{Physicist's summary.}
Beyond scaffolding, frame-change invariants select no realist position and
environment messages create no agents; coordinate ascent, learning, and delayed
feedback do not distinguish the program from hierarchical modeling, but
formation by the inference would.
